% Generated by run_retrospective_allocation_v2.py; verify cross-references before typesetting.
\subsection{Timestamp-controlled retrospective allocation replay}

Allocation validity was evaluated by replaying successfully delivered GetUsPPE batches using only acute-care requests timestamped before each decision. Donor availability was reconstructed from offers timestamped by the decision minus commitments timestamped earlier; the logged batch quantity was used only as a conditional dispatch budget and was excluded from evidence encoding. Lead risk used donor-to-recipient route strata rather than request age, and the cognitive operator was restricted and renormalized to the event-active graph while carrying each concept's previous converged state forward. Because the candidate count varied from tens to thousands, the softmax priority was divided by its active-set maximum before entering the LP; this outcome-independent normalization preserves its ordering and prevents the priority coefficient from vanishing as the candidate pool grows. The evidence encoder was calibrated with seeded TPE optimization on the calibration period, shortlisted configurations were selected on internal validation, the LP coefficient ratios remained fixed at the manuscript values, and test outcomes were not supplied to the optimizer. Relative to the frozen pre-optimization preset, tuning increased test recipient average precision from 0.073 to 0.081 and full-list NDCG from 0.222 to 0.233 (Holm-adjusted $p=0.0042$ for NDCG), while allocation overlap remained 0.008.

On the temporal test panels, coupled ACM obtained allocation overlap 0.008 [0.000, 0.016] and recipient average precision 0.081 [0.056, 0.161]. The strongest comparator on allocation overlap was Calibrated scalar at 0.110 [0.016, 0.202]. The Holm-adjusted paired comparison was not significant. These estimates measure agreement with historical recipient assignments conditional on batch size, not causal clinical benefit; donor preferences and shipment constraints absent from the public archive remain potential sources of disagreement.
